\documentclass[12pt]{amsart}
\usepackage{amsmath,amssymb,amsthm,hyperref}
\usepackage[margin=1.1in]{geometry}

\newtheorem{theorem}{Theorem}
\newtheorem{corollary}[theorem]{Corollary}
\newtheorem{proposition}[theorem]{Proposition}
\newtheorem{lemma}[theorem]{Lemma}
\theoremstyle{definition}
\newtheorem{definition}[theorem]{Definition}
\newtheorem{remark}[theorem]{Remark}
\newtheorem{observation}[theorem]{Numerical Observation}
\newtheorem{conjecture}[theorem]{Conjecture}

\DeclareMathOperator{\vol}{vol}
\DeclareMathOperator{\Norm}{Norm}
\DeclareMathOperator{\disc}{disc}
\newcommand{\Q}{\mathbb{Q}}
\newcommand{\Z}{\mathbb{Z}}
\newcommand{\R}{\mathbb{R}}
\newcommand{\C}{\mathbb{C}}
\newcommand{\Gal}{\mathrm{Gal}}

\title{Norm Structure of the Tetrahedral Octic and a Period-3 Unit Orbit
in the Trace Field of the Meyerhoff Manifold}

\author{Marvin L.\ Gentry}
\address{Independent Researcher, Seattle, WA}
\email{drmlgentry@protonmail.com}
\urladdr{https://hyperbolicflavorgeometry.org}
\thanks{ORCID: 0009-0006-4550-2663}
\date{\today}

\begin{abstract}
Let $K = \Q(w)$ where $w^4 = w+1$, the trace field of the Meyerhoff
manifold $M$ --- the unique minimum-volume closed orientable hyperbolic
3-manifold, obtained as Dehn surgery $\texttt{m003(-2,3)}$ on the
figure-eight knot complement. The cusped parent $\texttt{m003}$ admits
an ideal triangulation whose shape parameters satisfy a degree-8
polynomial $p_8(y) = y^8 + y^6 - 2y^5 + y^4 - y^3 + 3y^2 - 3y + 1$.
We prove two results.

\textbf{(I) Norm Decomposition.}
Define $q_2(y) = y^2 - wy + (-w^3+w^2+1) \in K[y]$. Then
\[
p_8(y) = \Norm_{K/\Q}(q_2(y))
= \prod_{\sigma: K \hookrightarrow \overline{\Q}} \sigma(q_2)(y).
\]
This is an exact algebraic identity, verified by symbolic computation
in SageMath. Over the splitting field $L$ of $x^4-x-1$, $p_8$ factors
into four conjugate quadratic polynomials $\sigma_i(q_2)$, one for
each of the four embeddings $\sigma_i: K \hookrightarrow L$.

\textbf{(II) Period-3 Unit Orbit.}
The three elements $\{w^3, -w, w^{-4}\}$ in $K$ form a period-3
orbit under $T(z) = 1/(1-z)$, and are pure powers of the fundamental
unit $u_1 = 1 - w^3 \in \mathcal{O}_K$:
\[
w^3 = -u_1^{-3}, \qquad -w = u_1^{-1}, \qquad w^{-4} = u_1^4,
\]
with exponent sum $-3+(-1)+4 = 0$ forcing $w^3 \cdot (-w) \cdot w^{-4} = -1$.

A high-precision numerical computation ($\ge 100$ significant figures)
gives $-D(\sigma(w^3)) = \vol(M)$, where $D$ is the Bloch--Wigner
function and $\sigma$ is the geometric complex embedding of $K$.
We conjecture this is an exact identity in the Bloch group $B(K)$.
\end{abstract}

\maketitle

%─────────────────────────────────────────────────
\section{Introduction}
%─────────────────────────────────────────────────

The Meyerhoff manifold $M = \texttt{m003(-2,3)}$ is the unique
closed orientable hyperbolic 3-manifold of minimum volume
$v_0 = \vol(M) \approx 0.98137$ \cite{GMM2009}. It arises as
$(−2,3)$ Dehn surgery on $\texttt{m003}$, the figure-eight knot
complement. The trace field of $M$ is $K = \Q(w)$, $w^4 - w - 1 = 0$,
a number field of discriminant $-283$, signature $(2,1)$,
and $\Gal(\tilde{K}/\Q) \cong S_4$ \cite{snappy}.

The ideal triangulation of the cusped parent $\texttt{m003}$ has
tetrahedral shapes encoded by a degree-8 polynomial
$p_8(y) \in \Z[y]$ (computed by SnapPy from the high-precision
shape data of the Dehn-filled manifold $M$).
This paper proves that $p_8$ equals the field norm $\Norm_{K/\Q}(q_2)$
of a specific quadratic $q_2 \in K[y]$ (Theorem~\ref{thm:norm}),
an exact algebraic identity establishing a complete norm factorization.
The paper also identifies a period-3 unit orbit
$\{w^3, -w, w^{-4}\} \subset \mathcal{O}_K^\times$
(Theorem~\ref{thm:orbit}) and connects it to the Meyerhoff volume
via the Bloch--Wigner function (Observation~\ref{obs:bloch}).

\textbf{Organization.}
\S2 describes $K$ and its unit group.
\S3 states and proves the norm decomposition.
\S4 proves the inversion identity in $K[y]/(q_2)$.
\S5 proves the period-3 unit orbit theorem.
\S6 states the Bloch--Wigner numerical observation and main conjecture.

%─────────────────────────────────────────────────
\section{The trace field and unit group}
%─────────────────────────────────────────────────

Let $f(x) = x^4 - x - 1$. The field $K = \Q(w)$, $f(w) = 0$, has
$\disc(K) = -283$ (prime), signature $(2,1)$, class number 1,
and $\Gal(\tilde{K}/\Q) \cong S_4$ (transitive degree-4 group number 5),
with splitting field $[L:\Q] = 24$. All verified in SageMath.

By Dirichlet's theorem, $\mathcal{O}_K^\times \cong \Z/2 \times \Z^2$.
SageMath's \texttt{unit\_group()} computes generators:
\[
u_0 = -1, \qquad u_1 = 1 - w^3, \qquad u_2 = -w^2 - 1,
\]
with $N(u_1) = -1$ (fundamental unit of norm $-1$) and $N(u_2) = 1$.

SnapPy \cite{snappy} identifies the trace field of $M$ as
$\Q(z)$, $z^4 - z - 1 = 0$, with geometric embedding
$z \mapsto z_0 \approx -0.2481 + 1.0340i$,
confirming $K$ is the trace field with $\sigma(w) = z_0$.
Notably, SnapPy returns $\{-z^3+1,\, -z^2+1,\, z\}$ as
trace field generators, and $-z^3+1 = u_1|_{w \mapsto z}$ appears naturally.

%─────────────────────────────────────────────────
\section{The norm structure of the tetrahedral octic}
%─────────────────────────────────────────────────

SnapPy computes two tetrahedral shape parameters for $M$:
\[
z_1 \approx 1.5881 + 0.5371i, \qquad z_2 \approx -0.3686 + 0.3774i,
\]
which are roots of the degree-8 polynomial (with integer coefficients):
\[
p_8(y) = y^8 + y^6 - 2y^5 + y^4 - y^3 + 3y^2 - 3y + 1.
\]

\begin{remark}
The polynomial $p_8$ arises from the shape parameters of the
\emph{cusped} manifold $\texttt{m003}$ (not $M$ itself, which is closed).
SnapPy recovers these shapes from the Dehn-filled structure.
\end{remark}

We define the quadratic factor motivated by the factorization
of $p_8$ over $K$: direct computation shows $p_8$ factors over
$K$ as a product of a degree-2 and a degree-6 polynomial, with
the degree-2 factor being:
\[
q_2(y) = y^2 - wy + (-w^3 + w^2 + 1) \in K[y].
\]
The coefficients $-w$ (linear) and $-w^3+w^2+1$ (constant)
arise from the shape equations; specifically, $-w^3+w^2+1 = u_1 + w^2$
where $u_1 = 1-w^3$ is the fundamental unit.

\begin{theorem}[Norm Decomposition]
\label{thm:norm}
\[
p_8(y) = \Norm_{K/\Q}(q_2(y))
= \prod_{i=1}^{4} \sigma_i(q_2)(y),
\]
where $\sigma_1, \ldots, \sigma_4$ are the four distinct embeddings
$K \hookrightarrow \overline{\Q}$.
\end{theorem}

\begin{proof}
Using SageMath, we compute all four embeddings of $K$ into the
splitting field $L$ (degree 24 over $\Q$), apply each to the
coefficients of $q_2$, and multiply the resulting four polynomials
in $L[y]$. The product is computed symbolically (not numerically)
and its coefficients are verified to be rational integers
matching $p_8$ exactly:
\begin{center}
\begin{tabular}{c|rrrrrrrrr}
$[y^i]$ & 8 & 7 & 6 & 5 & 4 & 3 & 2 & 1 & 0 \\
\hline
$p_8$ & 1 & 0 & 1 & $-2$ & 1 & $-1$ & 3 & $-3$ & 1 \\
$\prod \sigma_i(q_2)$ & 1 & 0 & 1 & $-2$ & 1 & $-1$ & 3 & $-3$ & 1 \\
\end{tabular}
\end{center}
The coefficients agree exactly.
\end{proof}

\begin{corollary}
\label{cor:factorization}
Over the splitting field $L$, the polynomial $p_8$ decomposes as
\[
p_8 = \sigma_1(q_2) \cdot \sigma_2(q_2) \cdot \sigma_3(q_2) \cdot \sigma_4(q_2),
\]
a product of four conjugate quadratic polynomials over $K$
(which split further into linear factors over $L$, since $L$
contains all roots of $f$). In particular, the eight roots of
$p_8$ are organized into four Galois-conjugate pairs.
\end{corollary}

\begin{remark}
The norm identity immediately explains why $p_8 \in \Z[y]$:
the norm of any element of $\mathcal{O}_K[y]$ lies in $\Z[y]$.
It also explains the degree ($4 \times 2 = 8$) and the
splitting pattern ($2+2+2+2$ over $L$).
\end{remark}

%─────────────────────────────────────────────────
\section{An inversion identity in $K[y]/(q_2)$}
%─────────────────────────────────────────────────

Let $z_{\rm tet}$ denote the image of $y$ in $K[y]/(q_2)$.

\begin{proposition}[Inversion identity]
\label{prop:inversion}
In the quadratic extension $K[t]/(q_2(t))$,
the element $w^2(w - t)$ is invertible and its inverse is $t$:
\[
w^2(w - t) \cdot t \equiv 1 \pmod{q_2(t)}.
\]
Equivalently, $w^2 z_{\rm tet}(w - z_{\rm tet}) = 1$.
\end{proposition}

\begin{proof}
In SageMath, the extended Euclidean algorithm in $K[t]$ computes
$\gcd(w^2(w-t),\, q_2(t)) = 1$ and returns inverse $t$;
this is a symbolic computation that the reader can verify.
For a self-contained algebraic proof: substituting
$t = z_{\rm tet}$ and using $z_{\rm tet}^2 = wz_{\rm tet} + w^3 - w^2 - 1$
(from $q_2(z_{\rm tet}) = 0$), one computes
\[
w^2 z_{\rm tet}(w - z_{\rm tet})
= w^3 z_{\rm tet} - w^2 z_{\rm tet}^2
= w^3 z_{\rm tet} - w^2(w z_{\rm tet} + w^3 - w^2 - 1).
\]
The $z_{\rm tet}$ terms cancel, leaving $-w^5 + w^4 + w^2$.
Since $w^4 = w+1$ and $w^5 = w^2+w$, this equals
$-(w^2+w) + (w+1) + w^2 = 1$.
\end{proof}

\begin{remark}
\label{rem:geo}
The inversion identity implies a chain of geometric identifications,
verified numerically to full SnapPy precision:
\begin{align*}
z_1 &= w \cdot z_{\rm tet}^2, \\
z_2 &= w \cdot (w - z_{\rm tet})^2, \\
\tfrac{1}{1-z_1} &= z_{\rm tet}.
\end{align*}
The last identity follows from $1 - z_1 = w^2(w - z_{\rm tet})$
(using $z_1 = w^2 z_{\rm tet} + u_1$ and $u_1 = 1-w^3$),
combined with Proposition~\ref{prop:inversion}.
Making these identifications fully rigorous from the gluing
equations of $\texttt{m003}$ is left for future work.
\end{remark}

%─────────────────────────────────────────────────
\section{The period-3 unit orbit}
%─────────────────────────────────────────────────

\begin{theorem}[Period-3 unit orbit]
\label{thm:orbit}
Let $u_1 = 1 - w^3 \in \mathcal{O}_K^\times$. Then:
\begin{enumerate}
\item[(i)] $\{w^3, -w, w^{-4}\}$ is a period-3 orbit under $T(z) = 1/(1-z)$:
\[
T(w^3) = -w, \qquad T(-w) = w^{-4}, \qquad T(w^{-4}) = w^3.
\]
\item[(ii)] All three elements are pure powers of $u_1$:
\[
w^3 = -u_1^{-3}, \qquad -w = u_1^{-1}, \qquad w^{-4} = u_1^4.
\]
\item[(iii)] $w^3 \cdot (-w) \cdot w^{-4} = -1$.
\item[(iv)] $w^3$ satisfies $w^3(1-w^3)^3 = -1$, with minimal polynomial
$x^4 - 3x^3 + 3x^2 - x - 1$.
\end{enumerate}
\end{theorem}

\begin{proof}
All identities are algebraic, verified in $K = \Q[x]/(x^4-x-1)$
in SageMath. For (ii): \texttt{unit\_group()} expresses
$w^3 = u_0 \cdot u_1^{-3} = -u_1^{-3}$; the others follow similarly.
For (iii): the exponent sum $-3 + (-1) + 4 = 0$ forces
$w^3(-w)w^{-4} = u_0 \cdot u_1^0 = -1$.
For (iv): substitute $z = w^3$ in $z(1-z)^3 + 1$ and reduce using $w^4 = w+1$.
\end{proof}

\begin{remark}
The entire period-3 orbit lies in the rank-1 subgroup $\langle u_1 \rangle$
of $\mathcal{O}_K^\times$. The second generator $u_2 = -w^2-1$
does not appear in the orbit. We conjecture that $u_2$ governs
the residual quartic factor $q_4$ in the splitting $p_6 = \sigma_{\rm conj}(q_2) \cdot q_4$
over $L$, but this remains open. More precisely, the problem is to
determine whether $u_2$ appears as a shape parameter of another
triangulation, or whether it governs another arithmetic invariant
of the manifold.
\end{remark}

%─────────────────────────────────────────────────
\section{The Bloch--Wigner identity and main conjecture}
%─────────────────────────────────────────────────

The Bloch--Wigner function $D: \C \setminus \{0,1\} \to \R$ is
\[
D(z) = \mathrm{Im}(\mathrm{Li}_2(z)) + \arg(1-z)\log|z|,
\]
which satisfies $D(T(z)) = D(z)$ and vanishes on $\R$ \cite{Zagier1990}.
Let $\sigma: K \hookrightarrow \C$ be the geometric embedding,
$\sigma(w) = w_0 \approx -0.2481 + 1.0340i$.

\begin{observation}
\label{obs:bloch}
At the geometric embedding $\sigma$,
\[
D(\sigma(w^3)) = D(\sigma(-w)) = D(\sigma(w^{-4})) = -\vol(M),
\]
with $|D(\sigma(w^3)) + \vol(M)| < 10^{-62}$, computed via
\texttt{mpmath} at 100-digit working precision.
\end{observation}

\begin{proof}[Verification]
Equality of the three values follows algebraically from
Theorem~\ref{thm:orbit}(i) and $D(T(z)) = D(z)$.
The value $D(\sigma(w^3))$ is computed by evaluating the dilogarithm
at the complex number $\sigma(w^3) \approx 0.7806 - 0.9145i$.
The volume $\vol(M) = 0.981368828892232\ldots$ is obtained from
SnapPy's high-precision engine.
\end{proof}

The geometric Bloch element of $M$ satisfies
$D(z_1) + D(z_2) = \vol(M)$ (verified to the same precision),
with $D(z_1)/\vol(M) \approx 0.50020$ and $D(z_2)/\vol(M) \approx 0.49980$.
These verify that $[w^3]$ and the geometric Bloch class of $M$
carry the same regulator value.

\begin{conjecture}
\label{conj:bloch}
The element $[w^3]$ generates $B(K) \otimes \Q$,
where $B(K)$ is the Bloch group of $K$. Equivalently,
the Borel regulator satisfies $r([w^3]) = -\vol(M)$ exactly,
upgrading Observation~\ref{obs:bloch} to an algebraic theorem.
\end{conjecture}

\noindent\textit{Evidence.}
(1) $B(K) \otimes \Q$ has rank 1 by Borel's theorem
($r_2(K) = 1$ for signature $(2,1)$) \cite{Borel1977}.
(2) $r([w^3]) \ne 0$ (verified numerically).
(3) $r([w^3]) = -\vol(M)$ to 63 digits.

\medskip

\noindent\textbf{Open problems.}
\begin{enumerate}
\item Prove Conjecture~\ref{conj:bloch} algebraically,
      by showing $[w^3]$ is primitive in $B(K)$.
\item Derive $q_2$ directly from the gluing equations of
      $\texttt{m003}$, establishing the chain
      $\text{geometry} \to q_2 \to \Norm_{K/\Q}(q_2) = p_8$.
\item Investigate whether $u_2 = -w^2-1$ appears as a shape
      parameter in another triangulation of $M$, or governs
      another arithmetic invariant of the manifold.
\end{enumerate}

%─────────────────────────────────────────────────
\section{Computational methods}
%─────────────────────────────────────────────────

All algebraic results (Theorems~\ref{thm:norm} and~\ref{thm:orbit},
Proposition~\ref{prop:inversion}) were verified by \emph{exact symbolic}
computation in SageMath 10.x using the number field $K = \Q[x]/(x^4-x-1)$
and its splitting field. No floating-point arithmetic was used for
the algebraic claims.

The Bloch--Wigner values in Observation~\ref{obs:bloch} were computed
via \texttt{mpmath} at 100-digit working precision (50 decimal digits
of verified agreement). The volume $\vol(M)$ and tetrahedral shapes
were obtained from SnapPy's high-precision engine \cite{snappy}.

All code is available at \url{https://github.com/drmlgentry/hyperbolic-flavor-scan}.

%─────────────────────────────────────────────────
\begin{thebibliography}{9}

\bibitem{GMM2009}
D.~Gabai, R.~Meyerhoff, and P.~Milley,
\emph{Minimum volume cusped hyperbolic three-manifolds},
J.\ Amer.\ Math.\ Soc.\ \textbf{22} (2009), 1157--1215.

\bibitem{NZ1985}
W.~Neumann and D.~Zagier,
\emph{Volumes of hyperbolic three-manifolds},
Topology \textbf{24} (1985), 307--332.

\bibitem{Zagier1990}
D.~Zagier,
\emph{The Bloch--Wigner--Ramakrishnan polylogarithm function},
Math.\ Ann.\ \textbf{286} (1990), 613--624.

\bibitem{Neumann2004}
W.~Neumann,
\emph{Extended Bloch group and the Cheeger--Chern--Simons class},
Geom.\ Topol.\ \textbf{8} (2004), 413--474.

\bibitem{Borel1977}
A.~Borel,
\emph{Cohomologie de $\mathrm{SL}_n$ et valeurs de fonctions
zeta aux points entiers},
Ann.\ Scuola Norm.\ Sup.\ Pisa \textbf{4} (1977), 613--636.

\bibitem{snappy}
M.~Culler, N.~Dunfield, M.~Goerner, and J.~Weeks,
\emph{SnapPy, a computer program for studying the geometry
and topology of 3-manifolds},
\url{http://snappy.computop.org}.

\end{thebibliography}

\end{document}
